\documentclass[a4paper,12pt]{article}
\usepackage[utf8]{inputenc}
\usepackage[T1]{fontenc}
\usepackage[ngerman]{babel}
\usepackage{geometry}
\geometry{margin=2.5cm}
\begin{document}

\section*{Deep Fluids: A Generative Network for Parameterized Fluid Simulations}

\textbf{Autoren:} Byungsoo Kim, Vinicius C.\ Azevedo, Nils Thuerey, Theodore Kim,
Markus Gross, Barbara Solenthaler \\
\textbf{Journal:} Computer Graphics Forum (Eurographics 2019), Vol.~38, Nr.~2 \\
\textbf{DOI:} 10.1111/cgf.13619

\subsection*{Zusammenfassung}

Kim et al.\ präsentieren ein generatives neuronales Netz, das in der Lage ist,
parametrisierte Fluidsimulationen direkt aus einem kompakten Satz von Steuerparametern
zu synthetisieren. Kernidee ist es, eine Kollektion von diskreten, parametrisierbaren
Strömungssimulationen zu nutzen, um ein faltendes neuronales Netz (CNN) zu trainieren,
das den zugrundeliegenden Parameterraum als nichtlinearen latenten Raum repräsentiert.

\paragraph{Architektur.}
Das Modell besteht aus einem generativen Netzwerk $G(\mathbf{c})$, das einen
Parametervektor $\mathbf{c} \in \mathbb{R}^n$ auf ein vollständiges Geschwindigkeitsfeld
$\hat{\mathbf{u}}_\mathbf{c}$ abbildet. Das Netzwerk beginnt mit vollständig verbundenen
Schichten, die die Eingangsparameter in einen kompakten Gewichtsvektor projizieren,
welcher anschließend durch eine Folge von Faltungsblöcken mit Upsampling-Operationen
und residualen Skip-Connections auf die volle Gitterauflösung expandiert wird. Für
erweiterte Parametrisierungen wird das Generatornetz durch einen symmetrischen
Encoder $G^\dagger(\mathbf{u})$ zu einem Autoencoder ergänzt, der beliebige
Geschwindigkeitsfelder in den latenten Raum $\mathbf{c} = [\mathbf{z}, \mathbf{p}]$
einbettet -- bestehend aus einem unüberwachten Teil~$\mathbf{z}$ und einem überwachten
Teil~$\mathbf{p}$.

\paragraph{Verlustfunktion und Divergenzfreiheit.}
Eine zentrale Neuerung der Arbeit ist eine strömungsfunktionsbasierte Verlustfunktion,
die Divergenzfreiheit der erzeugten Geschwindigkeitsfelder \emph{by construction}
garantiert:
\[
  L_G(\mathbf{c}) = \|\mathbf{u}_\mathbf{c} - \nabla \times G(\mathbf{c})\|_1.
\]
Da die Rotation eines beliebigen Vektorfelds stets divergenzfrei ist
($\nabla \cdot (\nabla \times G) = 0$), lernt das Netzwerk implizit eine
Stromfunktion~$\Psi_\mathbf{c}$. Zusätzlich wird ein Gradiententerm eingeführt,
der sicherstellt, dass auch die räumlichen Ableitungen des Geschwindigkeitsfeldes
korrekt approximiert werden -- ein Aspekt, der beim reinen $L_1$-Training im Nullraum
des Fehlermaßes verschwinden kann.

\paragraph{Latente Raumintegration.}
Für zeitlich veränderliche Szenarien, deren Parametrisierung über die Simulationsdauer
wächst, wird ein separates Integrationsnetzwerk $T(\mathbf{x}_t)$ eingeführt, das
Übergänge im latenten Raum erlernt und so eine Vorhersage neuer Zeitschritte direkt
im komprimierten Raum ermöglicht, ohne die vollständige Simulation erneut durchführen
zu müssen.

\paragraph{Ergebnisse.}
Das Verfahren wird an 2D- und 3D-Rauchsimulationen sowie an Flüssigkeitsszenarien
(Viskositätsvariationen, Kugelsedimentierung) demonstriert. Gegenüber dem zugrundeliegenden
CPU-Solver werden Beschleunigungen von bis zu $700\times$ erreicht, bei gleichzeitiger
Datenkompression von bis zu $1300\times$. Darüber hinaus zeigt das Modell plausible
Interpolationen für Parameterkombinationen, die nicht Teil des Trainingsdatensatzes waren.

\subsection*{Bezug zur Masterarbeit}

Die Arbeit von Kim et al.\ ist aus mehreren Perspektiven für die vorliegende Masterarbeit
relevant.

\paragraph{Gemeinsames Grundprinzip: Surrogatmodellierung für Strömungsfelder.}
Beide Arbeiten verfolgen das Ziel, rechenintensive numerische Strömungssimulationen
durch schnelle neuronale Netze zu ersetzen. Während Deep Fluids auf Navier-Stokes-Strömungen
mit inertialem Charakter (Rauch, Flüssigkeiten) fokussiert, adressiert die vorliegende
Masterarbeit Stokes-Strömungen im Niedrig-Reynolds-Regime, wie sie für magmatische Systeme
charakteristisch sind. In beiden Fällen ist das Trainingsziel ein physikalisch konsistentes
Geschwindigkeitsfeld auf einem regulären Gitter.

\paragraph{Divergenzfreiheit durch Stromfunktion.}
Die Kernidee von Deep Fluids, Divergenzfreiheit durch die Vorhersage einer Stromfunktion
$\Psi$ anstelle direkter Geschwindigkeitskomponenten zu erzwingen, ist ein zentrales
Strukturmerkmal der vorliegenden Masterarbeit. Das hier verwendete Lernziel ist
ebenfalls $\psi(x,z)$, die skalare Stromfunktion der zweidimensionalen Stokes-Gleichungen,
aus der die Geschwindigkeitskomponenten analytisch als $u_x = \partial_z\psi$ und
$u_z = -\partial_x\psi$ gewonnen werden. Diese Wahl der Ausgaberepräsentation
reduziert die effektiven Freiheitsgrade des Lernproblems und garantiert
Inkompressibilität ohne explizite Straffterme -- analog zu dem von Kim et al.\
motivierten Ansatz.

\paragraph{Unterschiede in Problemstellung und Architektur.}
Ein wesentlicher Unterschied besteht in der Problemstruktur: Deep Fluids lernt einen
\emph{parametrisierten} Lösungsraum, in dem Szenarien durch einen festen, niedrigdimensionalen
Parametervektor (z.\,B.\ Quellposition, Einstromgeschwindigkeit) beschrieben werden.
Im Gegensatz dazu ist das zentrale Generalisierungsproblem der vorliegenden Masterarbeit
die \emph{geometrische Variabilität}: Das Netz soll auf Konfigurationen mit variierender
Kristallanzahl und räumlicher Anordnung generalisieren -- also auf topologisch verschiedene
Eingaben. Diese Aufgabe ist nicht durch einen kompakten Parametervektor beschreibbar
und erfordert eine geometriegerechte Eingaberepräsentation (Kristallmaske, Distanzfeld,
Koordinatenkanäle) auf einem festen $256 \times 256$-Gitter.

Als Surrogatarchitekturen werden hier U-Net und Fourier Neural Operator (FNO) eingesetzt
und systematisch verglichen -- statt des in Deep Fluids genutzten generativen Decoder-Netzes
mit vollständig verbundenen Projektionsschichten. Der Fokus liegt nicht auf Interpolation
im latenten Raum, sondern auf der direkten Bild-zu-Bild-Abbildung von Geometrieinformation
auf das Stromfunktionsfeld.

\paragraph{Einordnung in die Literatur.}
Deep Fluids ist ein wichtiger Vorläufer datengetriebener Strömungssurrogate und motiviert
die physikalisch strukturierte Ausgaberepräsentation, die in der vorliegenden Arbeit
konsequent weiterverfolgt wird. Die hier untersuchte Frage der Generalisierung über
variable Partikelanzahlen -- also die Übertragung eines auf $N$ Kristallen trainierten
Modells auf Konfigurationen mit $N+m$ Kristallen -- ist in Deep Fluids nicht adressiert
und stellt einen eigenständigen Beitrag dieser Masterarbeit dar.

\end{document}